\documentclass[a4paper,fleqn]{cas-dc}

% ---------------------------------------------------------------------------
% Reporte técnico — GeoStats (plantilla Elsevier CAS, doble columna)
% Compilar:  cd docs/reporte && latexmk -pdf main.tex
% Cada sección vive en secciones/ y se incluye con \input (no \include:
% forzaría saltos de página y rompería las dos columnas).
% Las figuras vienen de los notebooks (notebooks/*.qmd) y viven en figuras/.
% ---------------------------------------------------------------------------

\usepackage[utf8]{inputenc}
\usepackage[T1]{fontenc}
\usepackage[spanish,es-tabla,es-noquoting,es-noshorthands]{babel}
\usepackage{amsmath,amssymb}
\usepackage{booktabs}
\usepackage{graphicx}
\graphicspath{{figuras/}}
\usepackage{multirow}
\usepackage{array}
\usepackage{xcolor}
\definecolor{azul}{HTML}{005991}
\definecolor{rojo}{HTML}{8B2C1A}
\definecolor{grafito}{HTML}{2C2C2C}
\hypersetup{colorlinks=true,linkcolor=azul,citecolor=azul,urlcolor=azul}

% Notación común
\newcommand{\NB}{\mathrm{BinomialNegativa}}
\newcommand{\Pois}{\mathrm{Poisson}}
\newcommand{\E}{\mathbb{E}}
\newcommand{\Var}{\mathrm{Var}}

% ---------------------------------------------------------------------------
% Localización de cas-dc al español. La clase trae en duro "ARTICLE INFO",
% "Preprint submitted to Elsevier", "Page x of y" y una nota "ORCID(s):"
% aunque no haya ninguno; se redefinen aquí sin tocar la clase.
% ---------------------------------------------------------------------------
\ExplSyntaxOn
\cs_set:Npn \__first_footerline:
  { \group_begin: \small \sffamily \__short_authors: \group_end: }
\cs_set:Npn \__first_foot:
  {
    \parbox[t]{\textwidth}
      { \rule{\textwidth}{.2pt} \\
        \__first_footerline: \hfill Página~\thepage{}~de~\lastpage }
  }
\cs_set:Npn \__cas_foot:
  {
    \parbox[t]{\textwidth}
      { \rule{\textwidth}{.2pt} \\ \sffamily \small
        \__first_footerline: \hfill Página~\thepage{}~de~\lastpage }
  }
\RenewDocumentCommand \printorcid { }
  {
    \seq_if_empty:NF \g_stm_orcid_seq
      {
        \group_begin:
        \tex_let:D \thefootnote \relax \footnotetext
          { \raggedright ORCID:~ \seq_use:Nn \g_stm_orcid_seq { ;~ } }
        \group_end:
      }
  }
\RenewDocumentEnvironment { keywords } { O{ Palabras~clave } }
  {
    \tex_global:D \tex_setbox:D \g_stm_key_box = \vtop \bgroup
    \hsize=.25 \textwidth
    \cs_set:Nn \sep: { \par }
    \cs_set_eq:NN \sep \sep:
    \parindent \c_zero_dim
    I\,N\,F\,O\,R\,M\,A\,C\,I\,Ó\,N \par \skip_vertical:n { -3pt }
    \rule{.25 \textwidth}{.2pt}\par\footnotesize
    \noindent \textit { #1 }:  \par
  }
  { \egroup }
\ExplSyntaxOff
% babel reimpone \abstractname al iniciar el documento; se fija en sus captions.
\addto\captionsspanish{\renewcommand{\abstractname}{R\,E\,S\,U\,M\,E\,N}}

% Encabezados cortos de página
\shorttitle{Predicción espacial de accidentes en la ZMM}
\shortauthors{E. Luna}

\begin{document}
\let\WriteBookmarks\relax
\def\floatpagepagefraction{1}
\def\textpagefraction{.001}

\title [mode = title]{Predicción espacial de accidentes de tránsito en la
  Zona Metropolitana de Monterrey}

\author[1]{Eduardo Luna}
\affiliation[1]{organization={Tecnológico de Monterrey},
  city={Monterrey}, state={Nuevo León}, country={México}}

% El entorno abstract de cas-dc es tipo verbatim: su \end{abstract} tiene que
% estar aquí; el texto vive en secciones/resumen.tex. El entorno descarta el
% primer token de su cuerpo: el \relax es ese sacrificio (sin él se pierde el
% \input). El "Overfull \hbox" que reporta \maketitle es propio de cas-dc.
\begin{abstract}
\relax% Texto del resumen (≈200 palabras): pregunta, datos, método y los tres
% hallazgos principales. Solo el texto: el entorno abstract vive en main.tex
% porque en cas-dc es tipo verbatim y no admite \end{abstract} en un \input.

\end{abstract}

\begin{keywords}
accidentes de tránsito \sep geoestadística \sep
modelo jerárquico bayesiano \sep CAR \sep Monterrey
\end{keywords}

\maketitle

\section{Introducción}\label{sec:introduccion}
% Motivación, la pregunta original (cuántos accidentes y dónde) y cómo se
% reformuló. Contribuciones del trabajo y organización del documento.

\section{Datos}\label{sec:datos}

\subsection{Fuente y recorte a la ZMM}\label{sec:datos-fuente}
\nocite{inegi_atus}% ATUS (INEGI), 2019--2024; los 18 municipios de la ZMM; llave (ANIO, EDO, MPIO, ID).

\subsection{Calidad y propiedades de la base}\label{sec:datos-calidad}
% Encoding, centinelas MNAR, panel nacional desbalanceado vs. ZMM balanceada.

\subsection{Selección de variables}\label{sec:datos-variables}

\section{Metodología}\label{sec:metodologia}

\subsection{Unidad espacial}\label{sec:met-unidad}
% Rejilla hexagonal en UTM 14N; por qué hexágonos y qué tamaño (MAUP).

\subsection{Análisis espacial exploratorio}\label{sec:met-exploratorio}
% Concentración (Lorenz), I de Moran, LISA y G* de Getis--Ord, persistencia.

\subsection{Modelos de conteo}\label{sec:met-conteo}
% Poisson, Binomial Negativa, tratamiento de ceros; panel y protocolo.

\subsection{Modelo jerárquico bayesiano}\label{sec:met-jerarquico}
% Especificación: efecto municipal dinámico + efecto espacial CAR por celda.

\subsection{Validación fuera de muestra}\label{sec:met-validacion}
% Entrenamiento 2019--2023, prueba 2024; métricas CRPS y PAI.

\section{Resultados}\label{sec:resultados}

\subsection{Estructura espacial de la siniestralidad}\label{sec:res-espacial}

\subsection{Modelos de conteo}\label{sec:res-conteo}

\subsection{Modelo jerárquico}\label{sec:res-jerarquico}
% Diagnósticos de ajuste y perfil espacial neto del municipio.

\subsection{Validación en 2024}\label{sec:res-validacion}

\section{Discusión}\label{sec:discusion}
% Régimen de reporte municipal como fuente dominante de error; lo que la base
% puede y no puede decir; limitaciones; caminos que no funcionaron.

\section{Conclusiones y trabajo futuro}\label{sec:conclusiones}


\appendix
\section{Supuestos del modelo jerárquico y su verificación}\label{app:supuestos}

\section{Reproducibilidad}\label{app:reproducibilidad}
% Comandos del pipeline (uv), notebooks Quarto, versiones de software, semilla.


\bibliographystyle{cas-model2-names}
\bibliography{referencias}

\end{document}
